\section{Corpus e separazione di TRAIN, DEV e TEST}\label{sec:stera30m-dati}

Il corpus di STERA 30M deriva dal file \texttt{TinyStoriesV2-GPT4-train} \cite{eldan2023tinystories}. I tre insiemi della \autoref{tab:stera30m-split} sono partizioni interne di quel file. Le storie sono tokenizzate con un vocabolario BPE di \num{50257} voci. L'identificativo 50.256 delimita l'inizio e la fine di ciascun racconto.

\begin{table}[htbp]
\centering
\small
\begin{tabularx}{\linewidth}{lrrX}
\toprule
Insieme & Racconti & Token target & Uso \\
\midrule
TRAIN & \num{2660865} & \num{530189191} & Aggiornamento dei parametri \\
DEV & \num{29279} & \num{5838521} & Controllo durante l'addestramento e valutazione su un campione distinto \\
TEST & \num{27351} & \num{5441693} & Prove predittive e narrative dopo il congelamento di STERA 30M \\
\bottomrule
\end{tabularx}
\caption{Partizioni del corpus di STERA 30M. I token target contano le posizioni predicibili di ogni racconto, compreso il token di fine.}
\label{tab:stera30m-split}
\end{table}

La separazione avviene per racconto. Il programma uniforma gli spazi bianchi e converte in minuscolo i byte ASCII prima di calcolare SHA-256. Interpreta i primi otto byte dell'impronta come intero in ordine \emph{little-endian} e ne prende il resto modulo \num{10000}. I valori da 0 a 99 vanno in TEST, quelli da 100 a 199 in DEV e gli altri in TRAIN. I \num{2000} gruppi di contenuto già riservati a una valutazione separata sono assegnati a DEV. Storie identiche dopo la normalizzazione ricevono la stessa assegnazione. L'audit ha verificato tutte le righe dell'indice. Testi simili ma con impronte diverse possono trovarsi in partizioni diverse.

\section{Configurazione del pretraining}
\label{sec:stera30m-config-training}

Il modello parte da zero ed esegue \num{122000} aggiornamenti, per un totale di \num{1529514422} token target: circa $2{,}9$ passaggi sull'insieme TRAIN. Ogni aggiornamento usa un batch di 64 sequenze, suddiviso fra due GPU NVIDIA P100, con accumulo dei gradienti su gruppi di otto sequenze per GPU. Le finestre arrivano a 512 token target e restano entro lo stesso racconto; il padding è escluso dalla perdita. L'ordine di visita di TRAIN è deterministico: il primo indice è \num{194723} e il passo è \num{104729} modulo il numero di racconti di TRAIN.

AdamW usa $\beta_1=0{,}9$, $\beta_2=0{,}95$ e limita la norma del gradiente a 1. Il weight decay è $0{,}1$ per coordinate RowMask23, fattori a basso rango e vettori di posizione. Il learning rate segue un profilo a tre fasi (\autoref{fig:stera30m-pretraining}): sale linearmente fino a $6\cdot10^{-4}$ nei primi 300 passi, resta costante fino al passo \num{64000}, scende linearmente fino a $6\cdot10^{-5}$ al passo \num{80000} e rimane a questo valore fino al termine. La NLL su DEV è misurata prima del training e ogni \num{2000} aggiornamenti, sempre sui primi 512 racconti DEV nell'ordine dell'indice (\num{102820} token target, al massimo 512 per racconto). Il modello valutato nei capitoli successivi è il checkpoint finale.

\begin{table}[htbp]
\centering
\small
\begin{tabularx}{\linewidth}{Xl}
\toprule
Operazione & Seme \\
\midrule
Seme globale PyTorch & \texttt{20261117} \\
Inizializzazione delle coordinate e dei vettori di posizione & \texttt{20261121} \\
Inizializzazione del fattore a basso rango $V$ & \texttt{7} \\
Selezione dei 300 incipit TEST & \texttt{202611240901} \\
Campionamento della continuazione $i$, con $i$ da 0 a 299 & $202611240902+1009\,i$ \\
\bottomrule
\end{tabularx}
\caption{Semi del pretraining e della valutazione. Il fattore $U$ parte da zero. Lo split è determinato dalla regola di hash della \autoref{sec:stera30m-dati}.}
\label{tab:stera30m-seed}
\end{table}
